% GENERATED FILE. Edit canonical lessons, then run npm run generate:books.
% canonical-source-hash: fnv1a64:d9661ce396010839
% canonical-lessons: KA-C479-sasya, KA-C479-pode, KA-C479-hullugavalu, KA-C479-ciguru, KA-C479-jaugu

\chapter{Grassland and Swamp}
\label{ch:ka-grassland-and-swamp}
\hlchaptermodality{479}

\begin{chapteropening}
\textbf{By the end of this chapter:} \emph{I can say vegetation, a bush, a grassland, to sprout and a swamp.}

Complete the last lesson of chapter 479: I can say vegetation, a bush, a grassland, to sprout and a swamp.
\end{chapteropening}

\section[\texorpdfstring{\kn{ಸಸ್ಯ} (sasya)}{ಸಸ್ಯ (sasya)}]{\kn{ಸಸ್ಯ} (sasya) — vegetation; a plant}

\label{lesson:KA-C479-sasya}



\begin{quote}
Before the new one: say the Kannada for a crab, then the Kannada for a whale.
\end{quote}

\subsection*{You'll want to know: \kn{ಸಸ್ಯ}}
\textbf{\kn{ಸಸ್ಯ}} — \emph{sasya} — \textquotedblleft{}vegetation; a plant\textquotedblright{}.

\begin{quote}
\kn{ಈ} \kn{ಬೆಟ್ಟದಲ್ಲಿ} \kn{ಅಪರೂಪದ} \kn{ಸಸ್ಯಗಳು} \kn{ಬೆಳೆಯುತ್ತವೆ}.
\end{quote}

\begin{quote}
\emph{ī beṭṭadalli aparūpada sasyagaḷu beḷeyuttave.}
\end{quote}

\begin{quote}
Rare plants grow on this hill.
\end{quote}

Plural: \kn{ಸಸ್ಯಗಳು}, \emph{sasyagaḷu}.

A word of science and botany; \kn{ಗಿಡ} is the everyday plant.

\begin{cousinweb}
\textbf{\kn{ಸಸ್ಯ}} (\emph{sasya}) is borrowed from Sanskrit \emph{sasya}, \textquotedblleft{}grain, fruit, crops\textquotedblright{} (CDIAL 13295); Kannada uses it for plants in general.
\end{cousinweb}

\begin{grammarlens}[title={What you've built}]
One more word for getting around.
\end{grammarlens}

\subsection*{Your turn}
\begin{itemize}
\raggedright
  \item \emph{Say it:} \emph{sasya}
  \item \emph{Say it:} \emph{sasya}, once more
  \item \emph{Say it:} \emph{ī beṭṭadalli aparūpada sasyagaḷu beḷeyuttave.} — Rare plants grow on this hill.
  \item \emph{Recall:} say the Kannada for a crab, then the Kannada for a whale, then say \emph{sasya} again
\end{itemize}

\begin{tcolorbox}[breakable,colback=teal!4,colframe=teal!35!black,title={Before you move on}]
What does \kn{ಸಸ್ಯ} mean? (\textquotedblleft{}Vegetation; a plant\textquotedblright{}.) Say it once more.
\end{tcolorbox}
\section[\texorpdfstring{\kn{ಪೊದೆ} (pode)}{ಪೊದೆ (pode)}]{\kn{ಪೊದೆ} (pode) — a bush; a thicket}

\label{lesson:KA-C479-pode}



\begin{quote}
Before the new one: say the Kannada for a whale, then the Kannada for vegetation.
\end{quote}

\subsection*{You'll want to know: \kn{ಪೊದೆ}}
\textbf{\kn{ಪೊದೆ}} — \emph{pode} — \textquotedblleft{}a bush; a thicket\textquotedblright{}.

\begin{quote}
\kn{ಮೊಲ} \kn{ಪೊದೆಯೊಳಗೆ} \kn{ಓಡಿತು}.
\end{quote}

\begin{quote}
\emph{mola podeyoḷage ōḍitu.}
\end{quote}

\begin{quote}
The rabbit ran into the bush.
\end{quote}

Low, thick growth where small animals hide.

\begin{cousinweb}
\textbf{\kn{ಪೊದೆ}} (\emph{pode}) is native Dravidian: DEDR 4509 lists Kannada \emph{pode}, \textquotedblleft{}a thatch, a bush, a bundle\textquotedblright{}, beside \emph{podaṟ}, \textquotedblleft{}a thicket\textquotedblright{}.
\end{cousinweb}

\begin{grammarlens}[title={What you've built}]
One more word for getting around.
\end{grammarlens}

\subsection*{Your turn}
\begin{itemize}
\raggedright
  \item \emph{Say it:} \emph{pode}
  \item \emph{Say it:} \emph{pode}, once more
  \item \emph{Say it:} \emph{mola podeyoḷage ōḍitu.} — The rabbit ran into the bush.
  \item \emph{Recall:} say the Kannada for a whale, then the Kannada for vegetation, then say \emph{pode} again
\end{itemize}

\begin{tcolorbox}[breakable,colback=teal!4,colframe=teal!35!black,title={Before you move on}]
What does \kn{ಪೊದೆ} mean? (\textquotedblleft{}A bush; a thicket\textquotedblright{}.) Say it once more.
\end{tcolorbox}
\section[\texorpdfstring{\kn{ಹುಲ್ಲುಗಾವಲು} (hullugāvalu)}{ಹುಲ್ಲುಗಾವಲು (hullugāvalu)}]{\kn{ಹುಲ್ಲುಗಾವಲು} (hullugāvalu) — a grassland; a pasture}

\label{lesson:KA-C479-hullugavalu}



\begin{quote}
Before the new one: say the Kannada for vegetation, then the Kannada for a bush.
\end{quote}

\subsection*{You'll want to know: \kn{ಹುಲ್ಲುಗಾವಲು}}
\textbf{\kn{ಹುಲ್ಲುಗಾವಲು}} — \emph{hullugāvalu} — \textquotedblleft{}a grassland; a pasture\textquotedblright{}.

\begin{quote}
\kn{ಹುಲ್ಲುಗಾವಲಿನಲ್ಲಿ} \kn{ಜಿಂಕೆಗಳು} \kn{ಮೇಯುತ್ತಿದ್ದವು}.
\end{quote}

\begin{quote}
\emph{hullugāvalinalli jiṅkegaḷu mēyuttiddavu.}
\end{quote}

\begin{quote}
Deer were grazing on the grassland.
\end{quote}

\kn{ಹುಲ್ಲು}, grass, plus \kn{ಕಾವಲು}, guarded land.

\begin{cousinweb}
\textbf{\kn{ಹುಲ್ಲುಗಾವಲು}} (\emph{hullugāvalu}) is a native Kannada compound of \kn{ಹುಲ್ಲು}, \textquotedblleft{}grass\textquotedblright{} (DEDR 4300), and \emph{kāvalu}, \textquotedblleft{}guarded land, a reserve\textquotedblright{}, built on \emph{kā}, \textquotedblleft{}to guard\textquotedblright{} (DEDR 1416).
\end{cousinweb}

\begin{grammarlens}[title={What you've built}]
One more word for getting around.
\end{grammarlens}

\subsection*{Your turn}
\begin{itemize}
\raggedright
  \item \emph{Say it:} \emph{hullugāvalu}
  \item \emph{Say it:} \emph{hullugāvalu}, once more
  \item \emph{Say it:} \emph{hullugāvalinalli jiṅkegaḷu mēyuttiddavu.} — Deer were grazing on the grassland.
  \item \emph{Recall:} say the Kannada for vegetation, then the Kannada for a bush, then say \emph{hullugāvalu} again
\end{itemize}

\begin{tcolorbox}[breakable,colback=teal!4,colframe=teal!35!black,title={Before you move on}]
What does \kn{ಹುಲ್ಲುಗಾವಲು} mean? (\textquotedblleft{}A grassland; a pasture\textquotedblright{}.) Say it once more.
\end{tcolorbox}
\section[\texorpdfstring{\kn{ಚಿಗುರು} (ciguru)}{ಚಿಗುರು (ciguru)}]{\kn{ಚಿಗುರು} (ciguru) — to sprout; to put out shoots}

\label{lesson:KA-C479-ciguru}



\begin{quote}
Before the new one: say the Kannada for a bush, then the Kannada for a grassland.
\end{quote}

\subsection*{You'll want to know: \kn{ಚಿಗುರು}}
\textbf{\kn{ಚಿಗುರು}} — \emph{ciguru} — \textquotedblleft{}to sprout; to put out shoots\textquotedblright{}.

\begin{quote}
\kn{ಮೊದಲ} \kn{ಮಳೆಯ} \kn{ನಂತರ} \kn{ಮರಗಳು} \kn{ಚಿಗುರಿದವು}.
\end{quote}

\begin{quote}
\emph{modala maḷeya nantara maragaḷu ciguridavu.}
\end{quote}

\begin{quote}
After the first rain the trees put out new shoots.
\end{quote}

Past (it …): \kn{ಚಿಗುರಿತು}, \emph{ciguritu}.

Also a noun, a tender shoot: \kn{ಹೊಸ} \kn{ಚಿಗುರು}.

\begin{cousinweb}
\textbf{\kn{ಚಿಗುರು}} (\emph{ciguru}) is native Dravidian: DEDR 2489 lists Kannada \emph{igaru}, \emph{cigur}, \emph{cigaru}, \textquotedblleft{}a sprout, a young leaf\textquotedblright{}, and \emph{cigi}, \textquotedblleft{}to sprout\textquotedblright{}.
\end{cousinweb}

\begin{grammarlens}[title={What you've built}]
One more word for getting around.
\end{grammarlens}

\subsection*{Your turn}
\begin{itemize}
\raggedright
  \item \emph{Say it:} \emph{ciguru}
  \item \emph{Say it:} \emph{ciguru}, once more
  \item \emph{Say it:} \emph{modala maḷeya nantara maragaḷu ciguridavu.} — After the first rain the trees put out new shoots.
  \item \emph{Recall:} say the Kannada for a bush, then the Kannada for a grassland, then say \emph{ciguru} again
\end{itemize}

\begin{tcolorbox}[breakable,colback=teal!4,colframe=teal!35!black,title={Before you move on}]
What does \kn{ಚಿಗುರು} mean? (\textquotedblleft{}To sprout; to put out shoots\textquotedblright{}.) Say it once more.
\end{tcolorbox}
\section[\texorpdfstring{\kn{ಜೌಗು} (jaugu)}{ಜೌಗು (jaugu)}]{\kn{ಜೌಗು} (jaugu) — a swamp; marshy ground}

\label{lesson:KA-C479-jaugu}



\begin{quote}
Before the new one: say the Kannada for a grassland, then the Kannada for to sprout.
\end{quote}

\subsection*{You'll want to know: \kn{ಜೌಗು}}
\textbf{\kn{ಜೌಗು}} — \emph{jaugu} — \textquotedblleft{}a swamp; marshy ground\textquotedblright{}.

\begin{quote}
\kn{ನದಿಯ} \kn{ಪಕ್ಕದ} \kn{ಜೌಗಿನಲ್ಲಿ} \kn{ಕಪ್ಪೆಗಳು} \kn{ಕೂಗುತ್ತಿದ್ದವು}.
\end{quote}

\begin{quote}
\emph{nadiya pakkada jauginalli kappegaḷu kūguttiddavu.}
\end{quote}

\begin{quote}
Frogs were calling in the marsh beside the river.
\end{quote}

\kn{ಜೌಗು} \kn{ನೆಲ}, marshy ground.

\begin{cousinweb}
\textbf{\kn{ಜೌಗು}} (\emph{jaugu}) is native Dravidian: DEDR 2398 lists Kannada \emph{javugu}, \emph{jōgu}, \textquotedblleft{}swampy ground\textquotedblright{}.
\end{cousinweb}

\begin{grammarlens}[title={What you've built}]
One more word for getting around.
\end{grammarlens}

\subsection*{Your turn}
\begin{itemize}
\raggedright
  \item \emph{Say it:} \emph{jaugu}
  \item \emph{Say it:} \emph{jaugu}, once more
  \item \emph{Say it:} \emph{nadiya pakkada jauginalli kappegaḷu kūguttiddavu.} — Frogs were calling in the marsh beside the river.
  \item \emph{Recall:} say the Kannada for a grassland, then the Kannada for to sprout, then say \emph{jaugu} again
\end{itemize}

\begin{tcolorbox}[breakable,colback=teal!4,colframe=teal!35!black,title={Before you move on}]
What does \kn{ಜೌಗು} mean? (\textquotedblleft{}A swamp; marshy ground\textquotedblright{}.) Say it once more.
\end{tcolorbox}
